\subsection*{一项旧请求可以挡住很多新结果}

取同一 E1 的序号 1、2、3。初始余额 10000，三项购买各为 1000 分，键分别为 K1、K2、K3。1 的请求迟迟没有确定响应，2、3 已完成，客户端已经持久收到它们的结果。连续确认水位为 0，服务暂时保留 2、3 的完整响应。这份保留来自连续范围编码的选择。

如果客户端继续发出 4、5 直到很大的序号，1 会阻塞后面全部响应回收。RIFL 为此限制单客户端不可回收请求的数量，论文给出其实现采用的 512 项上限~\cite{rifl}。本书可以选择未确认窗口 B，要求新序号落在 $h_{\rm ack}$+1 至 $h_{\rm ack}$+B 的范围内，超过窗口便先等待旧结果或处理空洞。

窗口只限制还没有被连续确认的范围。如果旧键映射留到 E1 关闭，已经确认的 n 份 $M_{\min}$/U 继续占空间。会话的单键数配额限制 n，固定期限约束它们保留多久。一个长期不结束的 C1 还会在会话关闭后继续占用恢复记录，系统要给这类任务安排独立配额、可追踪核查，并限制新增受理速度，避免新增任务长期超过终结速度。

\CTwelveFigure{ch12-original10}{序号 1 阻塞确认前缀，未确认窗口限制响应积压，会话配额另管逐键依据}{fig:ch12-window}

服务端看到序号 1 缺记录，也不能直接把水位改为 1。请求可能尚未到达，网络中的副本以后可进入。这样修改水位，会把一项本可首次受理的新请求错误标为旧请求。确认水位来自持久取得的终态响应；超时只表示一次通信没有给出这份响应。

\subsection*{跳过空洞要与迟到受理争同一个位置}

客户端可能在持久分配序号 1 和 S1 之后、第一次发送以前退出。恢复后确认自己希望放弃这项尚未受理的购买，可以请求关闭这个槽位。这里的槽位是某个协议客户端的一个序号，关闭槽位的命令记为 \code{SKIP(S1,1)}。它携带原 S1 和稳定编码，表示请求服务在尚未受理的情况下禁止它以后首次执行。

这项命令与 \code{ACCEPT(S1,1)} 在同一个受理域仲裁。SKIP 先胜出，就建立 U 的双向绑定和一条 \code{skipped} 身份记录。记录保存 S1、原 V/X 与稳定的跳过结果，P 字段为空，表示没有建立 P/C/T，并且未来同 S1 也不会被首次受理。后来迟到的同参 ACCEPT 读取它，返回原跳过结论，同 S1 异参则返回冲突。客户端持久取得这一终态以后，才可以把序号 1 纳入连续确认。

ACCEPT 先胜出时，SKIP 会读到已经受理的 P1/C1/T1。对于外部付款，响应是 \code{already\_accepted} 并给出原 P1/C1 的查询位置。C1 尚未终结，序号 1 继续未决，T1 沿已保存的平台身份核对原调用的结局。

\CTwelveFigure{ch12-original11}{SKIP 与 ACCEPT 同域仲裁，跳过先胜出禁止迟到受理，受理先胜出保留 C1/T1}{fig:ch12-skip}

只把放弃决定保存在客户端磁盘里，服务无法据此拦住已经在路上的请求。把取消命令发到另一张缺少共同唯一约束的表，也可能让取消与首次 A 分别提交。原键、原序号、跳过状态与首次受理使用同一状态转换，才有清楚的胜者。

这份跳过证据也进入 $M_{\min}$/U 的空间账。活动 E1 内，客户端换 q 再提交同一 K1，反向绑定继续拒绝。E1 关闭以后，整体资格接替逐笔跳过依据。若槽位尚未关联任何 S，可以单独封闭这个序号，但以后分配真实 K 时必须避开它。本书的具体流程从已经持久分配 S/q 的情况开始，保留一条可直接核对的链。

\subsection*{把确认与跳过接成完整状态转换}

下面以任意键外加序号的扩展路径给出活动 E1 的完整处理。可核验序号键从 K 解出 q，省去 U 的查找与写入；退休前缀统一拒绝旧序号，保留具体异参比较的部署继续查询原输入。$h_{\rm ack}$ 为已保存的确认水位，B 为未确认窗口上限，U 是双向映射。M/$M_{\min}$ 的槽位状态包括 \code{accepted}、已终结结果、\code{skipped} 与结果退休。\code{accepted} 在外部链路中指向已提交的 P/C/T，它的结果由原恢复任务补齐。

\begin{BookCode}
输入命令带完整 S、client_id、q、原输入及命令种类
以同一会话与身份仲裁域串行处理，先检查有效准入
若会话缺失、关闭或超期：返回 operation_expired
若时间依据不确定：返回 cutoff_uncertain
若 U(q) 指向别的 S，或 U(S) 指向别的 q：
    返回 encoding_conflict，保持原绑定
若已有 S 的绑定：
    按原 V/X 核对输入，异参返回 conflict
    若 q <= h_ack 或结果已退休：返回 result_retired
    若状态为 skipped：返回原 skipped 终态
    若为已完成且可回放：返回原 R(S)
    若为 accepted 且来件为 SKIP：
        返回 already_accepted 与原关联，保留 C/T 恢复
    若为 accepted 且来件为 ACCEPT：
        返回原受理信息和当前可核对进度
若尚无绑定且 q <= h_ack：
    报告 state_inconsistent，阻止首次业务并进入核查
若 q > h_ack+B：返回 confirmation_window_full
若尚无绑定且命令为 SKIP：
    共同保存 U 双向绑定、原 V/X 和稳定 skipped 结果
    该身份记录的 P 为空，不建立 P/C/T
    返回 skipped
若尚无绑定且命令为 ACCEPT：
    按既有购买规则仲裁
    共同保存 U、M 及 P/C/T 或本地终态
    返回实际保存的受理信息或终态
共同提交本次状态及应用边界以后响应
\end{BookCode}

已有身份分支先比较原 V/X，再判断结果是否退休。这样，确认水位 $h_{\rm ack}$ 以下的同 S 异参来件也会得到冲突，原绑定保持。若 q 已落入确认范围却完全找不到绑定，当前状态与已保存的确认事实相矛盾。服务通过 \code{state\_inconsistent} 停止这条首次受理路径，核对索引、快照与恢复来源。

\CTwelveFigure{ch12-slot-state}{空身份只由一次 SKIP 或 ACCEPT 离开，确认与退休保留逐键依据，未决受理沿原任务终结}{fig:ch12-slot-state}

ACK 单独使用同一个仲裁域推进连续范围：

\begin{BookCode}
处理 ACK(h_ack_new)：
    认证调用方，并核对所属 client_id 与 E1
    h_ack_new <= h_ack：确认已有水位，不删除其他状态
    对 h_ack+1 至 h_ack_new 的每个槽位核对已有终态
    缺槽位或仍有未决 accepted：
        拒绝推进，给出首个阻塞序号
    全部有终态：
        共同保存 h_ack_new，退休已确认响应
        活动 E1 的 M_min/U 继续保留
        C/T 按自身责任保存
    共同提交确认事实后响应
\end{BookCode}

ACK 的区间核对可以由终态范围索引实现，逻辑上覆盖每一个槽位。客户端的确认消息声明结果已经持久取得，服务端检查则排除把未完成槽位直接清掉的情况。ACK 先于结果终结到达，只会得到阻塞，客户端以后再次确认。

上述流程保留输入冲突、编码冲突和原结果各自的出口。已有 S 的支付输入来自原保存版本，新请求通过原绑定比较。确认和清理与业务写入若分开实现，就把可删除条件写进原子迁移：扫描提供候选，实际删除事务再次锁定并核对当前状态与确认事实，满足条件后才完成回收。

第十一章的主状态是 $(W_j,M_j,\Gamma_j,j)$。本节选用了 U 的扩展状态写为
\[
\Sigma^{\rm seq}_j=(W_j,M_j,\Gamma_j,U_j,j).
\]
其中 W 是业务数据，M 是逐笔绑定，$\Gamma$ 是会话准入状态，U 是可选的双向编码，j 是共同应用边界。本章的回收改变了其中各分量保存的内容，整份快照继续从同一 j 捕获它们。W 中的业务数据和已受理 C/T 保持原归属。M 按项保存完整结果或 $M_{\min}$，跳过也有可恢复的终态。$\Gamma$ 中每个活动会话保存账户、动作、固定截止、状态、确认水位和所需配额。U 保存编码的两个方向，分配进度随登记状态持久保存。

\CTwelveFigure{ch12-snapshot}{同一 j 的完整快照保存业务责任、缩减绑定、会话与确认事实，以及双向编码}{fig:ch12-snapshot}

现在可以检查同一 S 的受理次数。空初态中，每个 S 的次数都是零。第一次 ACCEPT 取得空身份时，原子建立唯一绑定，次数增加到一。之后同参恢复、异参或编码冲突、ACK 和响应退休都保留这份已见依据，次数保持。SKIP 先建立绑定时，次数为零，迟到 ACCEPT 从跳过记录退出。关闭以后，缺失或关闭 E1 的入口先拒绝，所以删除完成绑定也不会产生新的首次受理迁移。各允许迁移都保持每个 S 的受理次数至多一。

外部付款的效果上限再接上平台一侧：已经建立的 C 继续使用原范围与 \code{pay-C1}，所有可能生效的到达都有平台保护，才能把一次本地受理限制为至多一次外部成功。若平台保护随时间结束，就按下一节关闭或停止发送。活动 E1 提前删除最后一份逐键依据，会新增一条允许重复受理的迁移；保留 $M_{\min}$/U 或先关闭 E1，分别堵住这条迁移的身份入口和资格入口。

结果最终可回收还依赖进展条件。客户端能够恢复缺失结果、持久保存并发送确认，确认水位才会越过原空洞。客户端永久离线时，固定截止和可运行的关闭程序接替确认路径。未决 C/T 的终结则依赖恢复者获得执行机会、证据通道最终可用及本地 R 能提交。入口可以按期限停止接受旧请求，服务端未决责任的保存时间由证据决定。

\subsection*{关闭入口以后，未完成工作转到哪里}

现在让外部历史进入截止。Alice 的平台账户初始 10000，E1 在 10:00 截止。09:59，A 已提交 P1/C1/T1，金额为 3000，C1 为 \code{unconfirmed}，T1 为待恢复状态。平台请求可能未到达，也可能已经形成 X1、余额变为 7000。10:00 以后 E1 关闭。

购买入口的下一次旧请求返回 \code{operation\_expired}。恢复扫描器根据已受理 C1、原平台范围、\code{pay-C1} 和下游当前允许的继续方式工作。取得成功证据后，R 保存 C1 的两类原结果、O1、EV1，推进 P1，并结束 T1。取得足以终结的拒绝证据则保存原拒绝结果并结束任务。证据未知时，T1 继续出现在核查范围。

\CTwelveFigure{ch12-original12}{关闭 E1 终止旧购买入口，已受理 C1/T1 通过服务端恢复走到有证据的终态}{fig:ch12-close-recovery}

会话清理可以分段进行。先关闭并持久保存决定，保证所有入口都按权威状态判断。再等待关闭以前已经取得准入的 A 提交或回滚，收齐这批在途事务的归属。已完成且允许停止回放的逐笔记录可以删除。未完成 P/C/T 及其原 S、输入和平台关系保留。最后删除 E1 的活动登记行，缺失状态继续拒绝旧入口。

等待在途事务，在数据库里由登记锁与事务结束协调，在复制状态机里由已提交关闭位置和按序应用协调。外部发送线程又有自己的执行边界。第八章对资源端的要求是把“检查、水位推进和业务写入放在同一个原子决定中”。例如报告资源已经接受任务第 8 领取代次并保存水位 8，以后到达的第 7 领取代次覆盖写就会在该资源事务中被拒绝。T1 的领取到期只改变谁能提交本地进度，它不会撤回已经离开发送者的网络请求。

付款主线通过另一个接收边界合流。第八章给出的 W1、W2 即使分属任务第 7、8 领取代次，也都从原 C1 读取同一平台范围、\code{pay-C1} 与原付款输入。平台按这一身份保存首次受理与结果，两个到达共同找到 X1。若改为 \code{pay-C1-epoch8}，就会给平台另一项首次支付机会。因此 E1 关闭以后，C1 的恢复者继续使用固定支付身份，并逐次核对平台的保护窗口和安全发送条件。

\CTwelveFigure{ch12-fencing-boundaries}{本地领取、独立报告资源与平台付款各在自己的接收边界检查，支付身份沿 C1 保持}{fig:ch12-fencing-boundaries}

\subsection*{旧客户端重连时保留原来的未知}

Alice 的设备在 11:00 重启，只找到原 E1/K1 和确认输入，没有购买结果。客户端沿旧身份恢复，得到 410。这次响应准确告诉她旧入口已结束，但没有告诉她 X1 是否存在。她可以通过经账户授权的购买、支付事实查询找到 P1/C1。该查询沿已受理关联读取 \code{acceptance}、P1 当前 \code{progress} 和 C1 原 \code{payment\_result}，与购买准入分开。

若查询得到 P1 已付款，客户端保存原 P1/O1 和付款结果，结束旧任务。若查询得到原 C1 未决，界面展示核查中的购买，并使用对应查询位置等待证据。若当前查询没有找到记录，客户端保留原身份，结合查询来源、覆盖范围和旧请求可能到达的路径继续核查。

Ongaro 对会话缺失的处理可以放回这条历史理解。新客户端先调用 \code{RegisterClient} 登记，从集群取得新的会话编号。后续命令携带这个编号，状态机找不到对应会话时返回错误~\cite{ongaro}。\footnote{原文说明，遇到无会话记录的命令时，“it does not process the command and instead returns an error to the client”。其 LogCabin 示例进一步让该客户端退出，避免客户端误处理过期错误。本书把业务查询与旧购买入口分开，保留原身份供后续核查。}

假如 S1 已经完成，响应丢失，E1 后来被关闭并回收，自动重建 E1 会同时恢复它的首次受理资格，却已经找不到原逐笔结果。设备重连后携带的同一 K1 便可能再执行一次。另一条历史中，K1 从未到达，E1 也已关闭。两条历史在缺失会话入口上都返回过期，让关闭决定对整段旧输入保持相同含义。

\CTwelveFigure{ch12-missing-session}{缺失 E1 的两条历史统一拒绝旧提交，业务查询再沿可核验关联区分原购买事实}{fig:ch12-missing-session}

新的登记只产生从未使用的 E2。E2 下的购买须有独立业务决定，旧任务的未知结局继续沿 E1/K1 核查。TCP 重连、访问凭据刷新或服务进程更换只恢复通信与授权，客户端从持久任务中读回原业务意图、身份和输入，随后选择恢复原结果或查询原事实。
\FloatBarrier
